%% kapa: exam=1 round=29 year=2018 period=1 form=A subject=economics subject_ko=경제학원론 seq=30 number=70
%% source: Q-Net 감정평가사 기출문제 (cst00309) artlSeq=5202808 file=extracted/1008/5202808/2018년 제29회 감정평가사 1차시험 기출문제/2018년 제29회 감정평가사 1차시험 기출문제/1교시 A형.pdf
%% conversion: kordoc 4.24.51 -> tools/gen1.py (auto); review=proofread-sonnet; flags=none
\begin{question}{70}
경제성장이론에 관한 설명으로 옳은 것은?
\begin{choices}
  \item 내생적 성장이론(endogenous growth theory)에 따르면 저소득 국가는 고소득 국가보다 빨리 성장하여 수렴현상이 발생한다.
  \item 내생적 성장이론에 따르면 균제상태의 경제성장률은 외생적 기술진보 증가율이다.
  \item 솔로우 경제성장 모형에서 황금률은 경제성장률을 극대화하는 조건이다.
  \item 솔로우 경제성장 모형에서 인구 증가율이 감소하면, 균제상태에서의 1인당 소득은 감소한다.
  \item 솔로우 경제성장 모형에서 균제상태에 있으면, 총자본스톡 증가율과 인구 증가율이 같다.
\end{choices}
\end{question}
\answer{5}
